% Part:sets-functions-relations
% Chapter: size-of-sets
% Section: pairing-alt

\documentclass[../../../include/open-logic-section]{subfiles}

\begin{document}

\olfileid[ar]{sfr}{siz}{pai-alt}

\olsection{دالة مزاوجة بديلة}

\begin{explain}
ثمة تعدادات أخرى لـ$\Nat^2$ يسهل تعيين معكوساتها. وفيما يأتي
أحدها. فبدلًا من تمثيل التعداد بصريًا في مصفوفة، ابدأ بقائمة من
الأعداد الصحيحة الموجبة يقترن كل منها بموضع فارغ (في البداية). وتصور
ملء هذه المواضع تباعًا بالأزواج $\tuple{n,m}$ على النحو الآتي.
ابدأ بالأزواج التي لها~$0$ في المركبة الأولى (أي الأزواج
$\tuple{0,m}$)، وضع الأول منها (أي $\tuple{0,0}$) في أول موضع فارغ،
ثم تجاوز موضعًا فارغًا، وضع الثاني (أي $\tuple{0,2}$) في الموضع
الفارغ التالي، وتجاوز موضعًا مرة أخرى، وهكذا. وتبدو بداية تعدادنا
(غير المكتملة) الآن هكذا:
\[\small
\begin{array}{@{}c c c c c c c c c c c@{}}
\mathbf 1 & \mathbf 2 & \mathbf 3 & \mathbf 4 & \mathbf 5 & \mathbf 6 & \mathbf 7 & \mathbf 8 & \mathbf 9 & \mathbf{10} & \dots \\ \\
\tuple{0,0} &  & \tuple{0,1} &  & \tuple{0,2} &  & \tuple{0,3} & & \tuple{0,4} &  & \dots \\
\end{array}
\]
كرر ذلك بالأزواج $\tuple{1,m}$ في المواضع التي لا تزال فارغة،
متجاوزًا مرة أخرى موضعًا فارغًا من كل موضعين:
\[\small
\begin{array}{@{}c c c c c c c c c c c@{}}
\mathbf 1 & \mathbf 2 & \mathbf 3 & \mathbf 4 & \mathbf 5 & \mathbf 6 & \mathbf 7 & \mathbf 8 & \mathbf 9 & \mathbf{10} & \dots \\ \\
\tuple{0,0} & \tuple{1,0} & \tuple{0,1} &  & \tuple{0,2} & \tuple{1,1} &
\tuple{0,3} & & \tuple{0,4} &  \tuple{1,2} & \dots \\
\end{array}
\]
وأدخل الأزواج $\tuple{2,m}$، $\tuple{2,m}$، وهكذا، بالطريقة نفسها.
ومن ثم يبدأ تعدادنا المكتمل هكذا:
\[\small
\begin{array}{@{}cc c c c c c c c c c@{}}
\mathbf 1 & \mathbf 2 & \mathbf 3 & \mathbf 4 & \mathbf 5 & \mathbf 6 & \mathbf 7 & \mathbf 8 & \mathbf 9 & \mathbf{10} & \dots \\ \\
\tuple{0,0} & \tuple{1,0} & \tuple{0,1} & \tuple{2,0}  & \tuple{0,2} &
\tuple{1,1} & \tuple{0,3} & \tuple{3,0}  & \tuple{0,4} &  \tuple{1,2} & \dots \\
\end{array}
\]
وإذا رقّمنا خانات المصفوفة أعلاه وفقًا لهذا التعداد، فلن نجد خطًا
متعرجًا أنيقًا، بل الترتيب الآتي:
\[
\begin{array}{ c | c | c | c | c | c | c | c }
& \mathbf 0 & \mathbf 1 & \mathbf 2 & \mathbf 3 & \mathbf 4 & \mathbf 5 & \dots \\
\hline
\mathbf 0 & 1 & 3 & 5 & 7 & 9 & 11 & \dots \\
\hline
\mathbf 1 & 2 & 6 & 10 & 14 & 18 & \dots & \dots \\
\hline
\mathbf 2 & 4 & 12 & 20 & 28 & \dots & \dots & \dots \\
\hline
\mathbf 3 & 8 & 24 & 40 & \dots & \dots & \dots & \dots \\
\hline
\mathbf 4 & 16 & 48 & \dots & \dots & \dots & \dots & \dots \\
\hline
\mathbf 5 & 32 & \dots & \dots & \dots & \dots & \dots & \dots \\
\hline
\vdots & \vdots & \vdots & \vdots & \vdots & \vdots & \vdots & \ddots\\
\end{array}
\]
نرى أن الأزواج في الصف~$0$ تقع في المواضع ذات الأرقام الفردية في
تعدادنا، أي إن الزوج $\tuple{0,m}$ يقع في الموضع $2m+1$؛ وتقع أزواج
الصف الثاني، $\tuple{1,m}$، في المواضع التي تساوي أرقامها ضعف عدد
فردي، وتحديدًا $2 \cdot (2m+1)$؛ وتقع أزواج الصف الثالث،
$\tuple{2,m}$، في المواضع التي تساوي أرقامها أربعة أضعاف عدد فردي،
$4 \cdot (2m+1)$؛ وهكذا. وعوامل $(2m+1)$ لكل صف، $1$، $2$، $4$،
$8$، \dots، هي بالضبط قوى~$2$: أي $1= 2^0$، $2 = 2^1$،
$4 = 2^2$، $8 = 2^3$، \dots\@ والواقع أن الأس المعني هو دائمًا
العضو الأول من الزوج قيد النظر. ومن ثم يكون العامل للزوج
$\tuple{n,m}$ هو $2^n$. وهذا يعطينا الصيغة العامة:
$2^n \cdot (2m+1)$. لكن هذه دالة ترسل الأزواج إلى أعداد صحيحة
\emph{موجبة}، أي إن موضع $\tuple{0,0}$ هو~$1$. فإذا أردنا البدء من
الموضع~$0$، وجب أن نطرح~$1$ من الناتج. فنحصل على ما يأتي:
\end{explain}

\begin{ex}
الدالة $h\colon \Nat^2 \to \Nat$ المعطاة بـ
\[
h(n,m) = 2^n (2m+1) - 1
\]
هي دالة مزاوجة لمجموعة أزواج الأعداد الطبيعية~$\Nat^2$.
\end{ex}

\begin{explain}
وعليه، في تعدادنا الثاني لـ$\Nat^2$، رمز الزوج
$\tuple{0,0}$ هو $h(0,0) = 2^0(2\cdot 0+1) - 1 = 0$؛
ورمز $\tuple{1,2}$ هو $2^{1} \cdot (2 \cdot 2 + 1) - 1 = 2
\cdot 5 - 1 = 9$؛ ورمز $\tuple{2,6}$ هو $2^{2} \cdot (2
\cdot 6 + 1) - 1 = 51$.
\end{explain}

يكفي أحيانًا ترميز أزواج الأعداد الطبيعية~$\Nat^2$ من غير اشتراط
أن يكون الترميز شاملًا. ولمثل هذه الترميزات معكوسات ليست إلا دوالًا
جزئية.

\begin{ex}
الدالة $j\colon \Nat^2 \to \Nat^+$ المعطاة بـ
\[
j(n,m) = 2^n3^m
\]
هي دالة~!!a{injective} $\Nat^2 \to \Nat$.
\end{ex}

\end{document}
